\section{Splicing large content to low-complexity cells}
\label{sec:tpc31-all-content}

The preceding section assumes $c\le Z$.  TPC-30 supplies the
complementary estimate
\begin{equation}
 |\cS_\nu(c>C)|
 \ll_{\eps,h,\mathscr D}X^\eps
 \left(Q^2+\frac{XQ}{C}\right)
 \label{eq:tpc31-large-content-input}
\end{equation}
for every $C\ge1$ \citep{WangTPC30}.  Any further determinant
restriction only reduces the absolute majorant in that theorem.  We
now splice \eqref{eq:tpc31-large-content-input} to the small-content
cell estimate.

\subsection{A translated window in one residue class}

Let $I\subset\R$ be any interval of length at most $W$, and fix one
class $a\pmod q$.  Define the factorization-independent all-content
event
\begin{equation}
 \cE(I;a,q)
 :=\{\Delta^\#_{\alpha,\gamma}(j)\in I,
      \ \Delta^\#_{\alpha,\gamma}(j)\equiv a\!\pmod q\}.
 \label{eq:tpc31-all-content-cell-event}
\end{equation}

\begin{theorem}[All-content determinant-cell pruning]
\label{thm:tpc31-all-content-cell}
For every $1\le C\ll_{\mathscr D}Q$ and $1\le\nu\le3$,
\begin{align}
 |\cS_\nu(\cE(I;a,q))|
 \ll_{\eps,h,\mathscr D}{}&X^\eps
 \left\{
 Q^2+\frac{XQ}{C}
 \right.\notag\\[-0.2em]
 &\left.\quad
 +Q\log(2L)\left(1+\frac Wq\right)
 \bigl(C+J\log(2C)\bigr)
 \right\}
 \label{eq:tpc31-all-content-cell-expanded}\\
 \ll_{\eps,h,\mathscr D}{}&X^\eps XQ
 \left[
 \frac1J+\frac1C
 +\left(1+\frac Wq\right)
 \left(\frac CX+\frac1Q\right)
 \right].
 \label{eq:tpc31-all-content-cell}
\end{align}
The same estimate holds for the sum of the three raw channels.
\end{theorem}

\begin{proof}
Split the event according to $c>C$ and $c\le C$.  The first part is a
subevent of \eqref{eq:tpc31-large-content-input}.  For the second,
apply \cref{cor:tpc31-small-window-residue} with the same interval and
residue class for every $c\le C$.  This proves the expanded form.
Use $JQ\asymp X$ and absorb logarithms into $X^\eps$ to obtain
\eqref{eq:tpc31-all-content-cell}.
\end{proof}

The threshold $C$ is an optimization variable, not a condition in
the event.  We give the exponent consequence with fixed margins.

\begin{corollary}[Optimized all-content cell saving]
\label{cor:tpc31-optimized-cell}
Suppose
\begin{equation}
 \frac Wq+1\le X^{s+o(1)}
 \label{eq:tpc31-cell-capacity-exponent}
\end{equation}
for a fixed $s\ge0$.  If $s<\beta$, then every fixed
\begin{equation}
 \boxed{
 \sigma<\min\{1-\beta,\,\beta-s\}}
 \label{eq:tpc31-optimized-cell-saving}
\end{equation}
is a power-saving exponent relative to $XQ$ in
\eqref{eq:tpc31-all-content-cell}.
\end{corollary}

\begin{proof}
Fix a $\sigma$ satisfying
\eqref{eq:tpc31-optimized-cell-saving}.  Choose $\kappa$ strictly
between
\begin{equation}
 \sigma<\kappa<\min\{\beta,1-s-\sigma\}.
 \label{eq:tpc31-kappa-feasibility}
\end{equation}
This interval is nonempty: $\sigma<\beta$, and adding
$\sigma<1-\beta$ to $\sigma<\beta-s$ gives
\begin{equation}
 2\sigma<1-s.
 \label{eq:tpc31-kappa-average}
\end{equation}
Take $C=\lfloor X^\kappa\rfloor$.  The four ratios in
\eqref{eq:tpc31-all-content-cell} then save respectively
$1-\beta$, $\kappa$, $1-\kappa-s$, and $\beta-s$, each strictly more
than $\sigma$, up to $o(1)$.  The strict margins absorb the scale and
logarithmic losses.
\end{proof}

\begin{remark}[Endpoint convention]
\label{rem:tpc31-relaxed-endpoint}
Equation \eqref{eq:tpc31-optimized-cell-saving} is a relaxed margin:
every fixed exponent below the displayed endpoint is proved.  The
endpoint itself is not asserted, because the physical scales contain
$o(1)$ factors and the estimates contain logarithms.
\end{remark}

For an interval alone, take $q=1$.  If
$W=X^{\omega+o(1)}$ and $\omega<\beta$, the relaxed all-content
saving is
\begin{equation}
 \min\{1-\beta,\beta-\omega\}.
 \label{eq:tpc31-all-window-saving}
\end{equation}
The interval may be macroscopically separated from zero.

Such separation itself gives a further content collapse.

\begin{proposition}[Macro-centered content collapse]
\label{prop:tpc31-macro-content-collapse}
Fix $\lambda>0$, and suppose the translated interval $I$ lies in
\begin{equation}
 \{\delta\in\R:|\delta|\ge\lambda Q\}.
 \label{eq:tpc31-macro-center-condition}
\end{equation}
Then every physical triple with $\Delta^\#\in I$ has
\begin{equation}
 c_{\alpha,\gamma}(j)\ll_{\lambda,\mathscr D}1.
 \label{eq:tpc31-macro-content-bounded}
\end{equation}
Consequently, for an interval of length at most $W$ and one class
$a\pmod q$,
\begin{equation}
 \boxed{
 |\cS_\nu(\cE(I;a,q))|
 \ll_{\eps,h,\lambda,\mathscr D}
 X^\eps JQ\log(2L)\left(1+\frac Wq\right).}
 \label{eq:tpc31-macro-cell-bound}
\end{equation}
If $1+W/q\le X^{s+o(1)}$ with $s<\beta$, every fixed saving below
$\beta-s$ is available relative to $XQ$.
\end{proposition}

\begin{proof}
The physical row span gives
$|m_\alpha-m_\gamma|\le B_{\mathscr D}Q$.  Since
$m_\alpha-m_\gamma=c_{\alpha,\gamma}(j)\Delta^\#$,
\eqref{eq:tpc31-macro-center-condition} implies
$c_{\alpha,\gamma}(j)\le B_{\mathscr D}/\lambda$.  Apply
\cref{cor:tpc31-small-window-residue} with this fixed content ceiling.
The resulting content sum is $O_{\lambda,\mathscr D}(J)$, proving
\eqref{eq:tpc31-macro-cell-bound}.  Divide by
$XQ\asymp JQ^2$ for the exponent statement.
\end{proof}

\subsection{One residue class over the full range}

Let
\begin{equation}
 \cE(a,q)
 :=\{\Delta^\#_{\alpha,\gamma}(j)\equiv a\!\pmod q\}.
 \label{eq:tpc31-full-residue-event}
\end{equation}
Using the content-dependent range in
\eqref{eq:tpc31-full-residue-entropy} is sharper than enclosing every
content in one interval of length $O(Q)$.

\begin{theorem}[All-content residue-cell pruning]
\label{thm:tpc31-all-content-residue}
For every $1\le C\ll_{\mathscr D}Q$ and $1\le\nu\le3$,
\begin{align}
 |\cS_\nu(\cE(a,q))|
 \ll_{\eps,h,\mathscr D}{}&X^\eps
 \left(
 Q^2+\frac{XQ}{C}
 +QC+JQ\log(2C)
 \right.\notag\\[-0.2em]
 &\left.\hspace{8em}
 +\frac{Q^2}q\log(2C)+\frac{JQ^2}q
 \right)
 \label{eq:tpc31-all-residue-expanded}\\
 \ll_{\eps,h,\mathscr D}{}&X^\eps XQ
 \left(
 \frac1J+\frac1C+\frac CX+\frac1Q+\frac1q
 \right).
 \label{eq:tpc31-all-residue}
\end{align}
\end{theorem}

\begin{proof}
Use \eqref{eq:tpc31-large-content-input} for $c>C$ and
\cref{cor:tpc31-small-residue} with $Z=C$ for the complementary
part.  The simplified form follows as in
\cref{thm:tpc31-all-content-cell}.
\end{proof}

\begin{corollary}[Optimized residue-cell saving]
\label{cor:tpc31-optimized-residue}
If $q=X^{\rho+o(1)}$ with $\rho\ge0$, then every fixed
\begin{equation}
 \boxed{
 \sigma<\min\{1-\beta,\beta,\rho\}}
 \label{eq:tpc31-optimized-residue-saving}
\end{equation}
is a power-saving exponent relative to $XQ$, provided the minimum is
positive.
\end{corollary}

\begin{proof}
Fix a $\sigma$ satisfying
\eqref{eq:tpc31-optimized-residue-saving}.  Since
$\sigma<\min\{\beta,1-\beta\}$, choose
\begin{equation}
 \sigma<\kappa<\min\{\beta,1-\sigma\}.
 \label{eq:tpc31-residue-kappa-choice}
\end{equation}
Take $C=\lfloor X^\kappa\rfloor$.  In
\eqref{eq:tpc31-all-residue}, the five displayed ratios save
$1-\beta$, $\kappa$, $1-\kappa$, $\beta$, and $\rho$, respectively;
all exceed $\sigma$.  The strict margins absorb the endpoint losses.
\end{proof}

\begin{remark}[What the splice adds]
\label{rem:tpc31-splice-adds}
TPC-30 alone estimates the event $c>C$ without seeing determinant
support complexity.  The small-content theorem alone estimates a
determinant cell only after imposing $c\le C$.  Their splice proves a cellwise
statement for the entire retained raw packet.  It does not estimate
the union of all cells with a power saving.
\end{remark}
